\section{Implementacja}
\label{sec:implementation}

\subsection{Architektura rozwiązania}
\label{sec:implementation-architecture}
% Plan (etap 5): warstwa obliczeniowa (Rust) i warstwa prezentacji (Python)
% komunikujące się przez pliki; przepływ config -> train -> artefakty -> wykresy.

\subsection{Konfiguracja eksperymentu}
\label{sec:implementation-config}
% Plan (etap 5): plik config.toml (listing, tabela sekcji i kluczy), walidacja
% z diagnostyką wskazującą linię.

\subsection{Przygotowanie danych}
\label{sec:implementation-data}
% Plan (etap 5): wczytywanie CSV, kodowanie klas, normalizacja dopasowana do zbioru
% uczącego, podzbiór i wydzielenie walidacji.

\subsection{Sieć i uczenie}
\label{sec:implementation-network}
% Plan (etap 5): warstwy jako macierze ndarray, asercje wymiarów, pętla uczenia,
% obserwator, wykrywanie dywergencji, opcjonalny early stopping.

\subsection{Powtarzalność}
\label{sec:implementation-reproducibility}
% Plan (etap 5): jedno ziarno, osobne strumienie RNG (dane, inicjalizacja, tasowanie,
% próbkowanie wag), wersja kodu w metadata.json.

\subsection{Weryfikacja poprawności}
\label{sec:implementation-verification}
% Plan (etap 5): numeryczny gradient check (różnice centralne), testy determinizmu,
% CI.

\subsection{Artefakty, wizualizacja i automatyzacja badań}
\label{sec:implementation-analysis}
% Plan (etap 5): artefakty (.npy, epochs.csv, metadata.json), wykresy w Pythonie,
% study.toml -> generate -> train -> aggregate.
